%% ============================================================================================
%% DRAFT v2, part 4: Discussion and implications.
%% ============================================================================================

\section{Discussion}
\label{sec:discussion}

\subsection{Withholding permission is not withholding authority}
\label{sec:disc-leastpriv}

Specifications meant to inspect are given file-writing tools far less often than those meant to change code, and the
gap survives every curation cut and is paired within repositories. Whether written or copied, these grants withhold
\code{Write} from reviewers. But the withholding applies to permissions, and a tool allowlist enumerates permissions,
not authority. Because a shell carries the process's ambient file-system authority \citep{miller2003capability},
removing the file-writing tools $W=\{\code{Write},\code{Edit},\code{NotebookEdit}\}$ from a grant $T$ that retains an
unsandboxed \code{Bash} leaves its authority over files unchanged, $A(T \setminus W) = A(T)$, whenever shell commands run without per-command
approval. In that context the set $\{\code{Read}, \code{Grep}, \code{Glob}, \code{Bash}\}$ looks like a reader and acts
like a read--write agent; in \code{default} mode the user's approval of each command stands between the two.

A retained shell is not always a mistake. Some write-withholding specifications ask for builds and tests, and a
minority ask for state-changing commands, but most need at most read-only commands or none (\S\ref{sec:res-rq3}). For those, the shell adds authority the text never asks for; for the rest, the problem is that
the tool list offers no way to grant part of a shell's reach.

\paragraph{Why the abstraction invites the error.}
\emph{Objects at different levels.} Lampson's access matrix treats each object as an independent column whose rights
are granted separately \citep{lampson1974protection}; the \code{tools} field mixes resource operations (\code{Read},
\code{Write}, \code{Edit}) with an interpreter (\code{Bash}) that can perform any of them.
\emph{Direct paths are guarded, indirect paths are not.} Tools tend to guard the obvious interface and miss indirect
routes to the same resource \citep[\S4.2]{garfinkel2003traps}; file-editing deny rules do not cover subprocesses that
write files indirectly \citep{AnthropicPermissionsDoc}.
\emph{Fail-safe defaults are inverted.} Omitting \code{tools} grants everything, so the least deliberate configuration
is the most privileged one, contrary to basing access on permission rather than exclusion
\citep{saltzer1975protection}; \RQtwoImplicitEng{} of specifications omit the field.
\emph{The resolver steers delegates to the shell.} It removes \code{Grep} and \code{Glob} whenever \code{Bash}
resolves, which moves routine search onto the shell, and leaves \code{Write} and \code{Edit} in place without a
warning.

\subsection{An agent-configuration security smell}
\label{sec:disc-smell}

Coarse settings that defeat finer controls beside them recur in configuration security (\S\ref{sec:rw-smells});
Table~\ref{tab:smells} places our pattern in that family. We name \textbf{unscoped shell in a restricted role} as a candidate security smell: a
specification that withholds file-writing tools while granting a shell without a command scope. It maps to CWE-250
(execution with unnecessary privileges) and to CWE-653 (improper isolation or compartmentalisation), because the file
tools were separated from the grant while one tool that spans them was not \citep{cwe250,cwe653}. It matters only
where shell commands run without per-command approval and without a sandbox (\S\ref{sec:capmodel}). The linter
released with this paper reports the specification-level pattern; it does not consult permission modes, settings or
hooks, and neither the smell nor the linter's precision has been validated with practitioners.

\begin{table}[t]
\caption{Coarse or over-broad grants that defeat least privilege or finer controls beside them. The last row is the
pattern measured in this paper.}
\label{tab:smells}
\footnotesize
\begin{tabular}{@{}>{\raggedright\arraybackslash}p{3.3cm}>{\raggedright\arraybackslash}p{5.1cm}>{\raggedright\arraybackslash}p{2.5cm}@{}}
\toprule
\textbf{Coarse grant} & \textbf{What it defeats} & \textbf{Source} \\
\midrule
\code{privileged: true} in a Kubernetes security context & every other access control in the same security context & \citet{rahman2023security} \\
default read--write \code{GITHUB\_TOKEN} & the read-only needs of nearly all workflows & \citet{koishybayev2022characterizing} \\
admin account as the default user & a default least-privileged account (CWE-250) & \citet{rahman2019seven} \\
\code{pledge} with \code{exec} and no \code{execpromises} & every promise, for the child process & \citet{openbsd_pledge} \\
\textbf{unscoped \code{Bash} in a subagent tool list} & \textbf{the withholding of \code{Write}, \code{Edit} and \code{NotebookEdit} in the same list, when shell commands run without approval and without a sandbox} & \textbf{this paper} \\
\bottomrule
\end{tabular}
\end{table}


\subsection{Prompts are interpreted, grants are enforced}
\label{sec:disc-prompts}

That an agent with a shell can write when asked and commands are approved automatically is expected, and the
mechanism check confirms it. Three results matter more for practice. First, in the realistic inspection configuration
the prompt restriction worked for Sonnet and Opus and left Haiku completing \ExpCthreeHaiku{} of the changes it
forbids. A restriction that holds or fails with the
model is a property of the model, not of the configuration, and \RQfiveNoModelPooled{} of specifications (pooled;
repository mean \RQfiveNoModel{}) do not name the model that will read them. Second, the channels fail differently: a
missing tool produces an error the parent can see, an unheeded prompt a completed change, and Haiku's false claims
show that the delegate's report is not a reliable record either. Third, on one fixture with explicitly report-only tasks, no agent
changed files under any grant. For such tasks, the risk of a dominated withholding lies less in delegates acting
unasked than in a grant that does not refuse a change when one is requested. Our requests came from the parent;
whether content the delegate reads can produce such a request \citep{debenedetti2024agentdojo} we did not test. The
enforcement-channel fields of a specification (\code{tools}, \code{disallowedTools}, \code{permissionMode},
\code{hooks}, \code{isolation}) are the only ones that do not depend on the model agreeing, and a prompt restriction
that holds for some models is not a security control.

Two trust boundaries lie outside our measurements. A delegate's answer returns to a parent that usually holds more
authority, so restricting the delegate does not stop content it reads from reaching the parent; Haiku's false
completion claims are a benign case of a parent receiving unreliable output. And a delegate that can write files, by
a tool or through a shell, can edit agent specifications, settings or git hooks and so widen grants in later
sessions.

\citet{Yan2026DoNotDeny} found that few security rules in context files have an enforced counterpart; we find the
same gap inside one artifact: when a specification's text makes the agent read-only, its grant can usually still reach
the file system (\RQfourFullReaches{}).

\subsection{Specifications as maintained artifacts}
\label{sec:disc-drift}

Most files naming a removed tool were committed after the removal (\S\ref{sec:res-rq6}); the names disappeared between
\DriftRemovalWindow{}, and \RQsixDriftInSnapshotYear{} of these files were first committed in 2026. Stale names thus
enter at authoring time, which is consistent with the reuse we measure and with copy-based reuse of skills and editor
rules \citep{Destefanis2026GitSkills,Jiang2026CursorRules}. Revalidation alone does not address this; a resolver that
reports what it dropped does.

\subsection{Implications}
\label{sec:disc-implications}

\paragraph{For practitioners.} Omitting \code{tools} grants everything. An agent that holds an unscoped shell is not
read-only, whatever its prompt says, unless each shell command needs approval or runs in a sandbox. Where the text asks
only to read, restrict the shell with Bash deny rules in settings, which the documentation says apply to
subagents \citep{AnthropicSubagentsDoc}, but such rules match command text and can be bypassed \citep{Chen2026Denylist}. A command pattern in
a subagent's \code{tools} field is not a restriction: in our test it granted the whole shell (\S\ref{sec:res-rq3}).
Where the restriction must hold, or the text asks to build or
test, use a sandbox or a writable scratch directory. Do not treat a prompt restriction as a security control; if one matters, pin the model and
test it. Check that tool names resolve on the release you run and that every \code{name} in an agent directory is
unique. The linter released with this paper reports all of these and exits non-zero on dominated withholdings,
ignored files and duplicate names.

\paragraph{For tool builders.} \emph{Report what was dropped}: the settings layer warns about a deny rule that matches
no tool, while the frontmatter resolver says nothing about a discarded name, a name collision or a file it does not
load. \emph{Warn on dominated withholding}: a specification that withholds file-writing tools and grants an unscoped
shell may not be what its author intended. \emph{Make read-only an execution-environment property}, enforced where every subprocess is covered, by a sandbox
or capability mode \citep{watson2010capsicum}. \emph{Make grants typed and attenuable}: a shell grant should carry a
file-system and network scope \citep{moore2014shill}, and policies should inspect arguments rather than tool names
\citep{shi2025progent}, outside the model \citep{debenedetti2025defeating}. Where a narrow tool
can replace the shell, ship the narrow tool \citep{beurerkellner2025design}.

\paragraph{For researchers.} Measure the effective configuration by executing its consumer, and test behavioural
claims by repeated execution, including tasks that do not ask for the behaviour.
